%--------------------------------------------------------------------------
%		Não é necessário mudar qualquer coisa no preambulo deste arquivo
%		basta apenas colocar o seu texto, e as palavras-chave
%		observando as normas ABNT
%--------------------------------------------------------------------------
\chapter*{RESUMO} 
\thispagestyle{empty}
O resumo é um elemento obrigatório, e segundo a norma NBR6028:2003, ``Consiste na apresentação concisa do trabalho'', \cite[p.~1]{NBR6028:2003}. Deve descrever de forma clara e sintética a natureza do trabalho, o objetivo, o método, os resultados e as conclusões, visando fornecer elementos para o leitor decidir sobre a consulta do trabalho no todo. O resumo deve destacar: o assunto do trabalho; o objetivo do texto; as conclusões do trabalho. Deve ser redigido em linguagem clara e objetiva; ser inteligível por si mesmo (dispensar a consulta ao trabalho), evitar repetição de frases inteiras do trabalho, respeitar a ordem em que as ideias ou fatos são apresentados, empregar verbos na voz passiva, utilizar de preferência a 3ª pessoa do singular, evitar o uso de frases negativas, abreviaturas, fórmulas, quadros, equações, etc.
Para teses e dissertações, monografias, relatórios e TCCs o resumo deve ser
de 150 a 500 palavras. Deve ser redigido em parágrafo único, com frases completas,
sem sequência de itens/tópicos. Seguidas do resumo, as palavras-chave e descritores devem ser representativas ao conteúdo do trabalho, sendo as palavras-chave separadas entre si por ponto e finalizadas também por ponto. 
 
 
%--------------------------------------------------------------------------
\vspace{2cm}
{\Large \textbf{Palavras-chave:}}
\vspace{0.4cm}
%--------------------------------------------------------------------------

Latex. Modelos em latex. Monografia. 
%MAIS DICAS
%
%Passo 1: O que é (contexto)
%
%Para iniciar o resumo e abstract, apresente ao leitor do que se trata seu trabalho, ou seja, qual problema você está investigando e a sua relevância.
%
%Logo aqui, nas primeiras linhas, você precisa criar um contexto que desperte a curiosidade e o interesse para a leitura.
%
%Passo 2: Por quê? (Descreva o objetivo do trabalho)
%
%Com base no problema ou situação investigada, seu trabalho já precisa dizer a que veio. A forma de descrever um objetivo varia de acordo com o estudo realizado.
%
%Seu trabalho pode se propor a analisar diferentes aspectos de um mesmo problema, a compreender um fenômeno, ou então, a “apresentar uma solução para reduzir a produção de energia não renovável”, por exemplo. Em resumo, você conta o porquê do seu TCC existir.
%
%Passo 3: Como? (Método utilizado)
%
%Você deve explicar rapidamente como seu trabalho foi feito, qual metodologia utilizou. Conte ao seu leitor se você realizou um estudo de caso, um experimento, uma pesquisa quantitativa e/ou pesquisa qualitativa etc.
%
%Passo 4: Resultados
%
%Apresente os resultados mais relevantes do seu estudo, sejam eles positivos ou negativos. Descreva-os de forma sucinta, e que despertem curiosidade e interesse em quem lê.
%
%Passo 5: Conclui-se que…
%
%Por fim, seu resumo deve trazer as principais conclusões sobre tudo o que você leu, pesquisou e desenvolveu.
%
%Aqui, também é possível mencionar como o seu trabalho poderá contribuir para o entendimento do objeto ou fenômeno estudado, ou como servirá de base para pesquisas futuras.
%\vspace{-1cm}


